\documentclass[a4paper,12pt]{article}
\usepackage{docsTemplate}
\usepackage{usecase}
\usepackage{float}

\setTitle{Analisi dei Requisiti}
\setAuthors{Filippo Compagno \\ & Sara Gioia Fichera \\ & Andrea Masiero}
\setVerificators{Andrea Masiero \\ & Francesco Balestro \\ & Bianca Zaghetto \\ & Sara Gioia Fichera}
\setApprovation{}
\setVersion{0.5}
\setType{Esterno}
\setDestination{Tullio Vardanega \\ & Riccardo Cardin \\ & Sanmarco Informatica \\[-0.1em] & Team Three Technologies}

% -- Obbligatori --
\newcommand{\loaddip}{01}
\newcommand{\remloaddip}{02}
\newcommand{\errloaddip}{03}
\newcommand{\removedip}{04}
\newcommand{\viewdip}{05}
\newcommand{\viewdetdoc}{06}
\newcommand{\viewconsinfo}{07}

% IdDoc
\newcommand{\mdparamImpCrit}{08} % Impronta crittografica - Hash - Algoritmo
\newcommand{\mdparamIdDoc}{09} % Identificativo documento
% ModalitàDiFormazione
\newcommand{\mdparamModFor}{10} % Modalità di formazione
% TipologiaDocumentale
\newcommand{\mdparamTipDoc}{11} % Tipologia documentale
% DatiDiRegistrazione
\newcommand{\mdparamTipFlu}{12} % Tipologia di flusso
\newcommand{\mdparamTipReg}{13} % Tipo registro
\newcommand{\mdparamDatReg}{14} % Data registrazione
\newcommand{\mdparamOraReg}{15} % Ora registrazione
\newcommand{\mdparamNumDoc}{16} % Numero documento
\newcommand{\mdparamCodReg}{17} % Codice registro

\begin{document}
    \include{docTitlepage}

    \section*{Registro delle modifiche} {
        \begin{table}[h!]
            \rowcolors{2}{lightgray}{white}
                \begin{tabularx}{\textwidth}{|l|l|l|l|L|L|}
                \hline
                \textbf{Vers.} & \textbf{Data} & \textbf{Autore} & \textbf{Ruolo} & \textbf{Verifica} & \textbf{Oggetto} \\
                \hline
                \textbf{0.5} & 04/12/25 & Andrea Masiero & Analista & Sara Gioia Fichera & Inserimento Casi d'Uso \\
                \hline
                \textbf{0.4} & 01/12/25 & Filippo Compagno & Analista & Bianca Zaghetto & Inserimento Casi d'Uso \\
                \hline
                \textbf{0.3} & 24/11/25 & Sara Gioia Fichera & Analista & Francesco Balestro &  Definizione Attori \\
                \hline
                \textbf{0.2} & 20/11/25 & Filippo Compagno & Analista & Francesco Balestro & Caratteristiche degli utenti \\
                \hline
                \textbf{0.1} & 07/11/25 & Sara Gioia Fichera & Analista & Andrea Masiero & Struttura iniziale del documento \\
                \hline
            \end{tabularx}
        \end{table}
    }
    \newpage
    
    \tableofcontents
    \newpage

    \listoffigures
    \newpage
    
    \listoftables
    \newpage

    \section{Introduzione} {
        \subsection{Scopo del documento} {        
            Questo documento ha lo scopo di descrivere in maniera esaustivo e precisa i requisiti funzionali e non funzionali del sistema software da sviluppare per il progetto DIPReader. Per questo motivo integra un'attenta descizione dei casi d'uso, che sono fonte di requisiti.
            Tale analisi è il risultato di un attento studio del capitolato, integrato con quanto concordato durante gli incontri svolti con Sanmarco Informatica, la proponente.
        }
        \subsection{Glossario} {
            Si rimanda al documento Glossario per chiarire i termini nel dominio del progetto che possono risultare ambigui o di difficile comprensione. I vocaboli presenti nel glossario sono stati segnalati seguendo la notazione:
            \begin{center}
                parola\textsubscript{G}
            \end{center}
        }
        \subsection{Riferimenti} {
            \subsubsection{Riferimenti normativi} {
                \begin{itemize}
                    \item Norme di Progetto
                    \item \hypertarget{capitolato}{Capitolato}
                    \begin{itemize}
                        \item {\url{https://www.math.unipd.it/~tullio/IS-1/2025/Progetto/C3.pdf}}
                        \item[] (Ultimo accesso: 2025/11/07)
                    \end{itemize}
                    \item Regolamento del Progetto Didattico
                    \begin{itemize}
                        \item \url{https://www.math.unipd.it/~tullio/IS-1/2025/Dispense/PD1.pdf}
                        \item[] (Ultimo accesso: 2025/11/07)
                    \end{itemize}
                    \item 830-1998 - IEEE Recommended Practice for Software Requirements Specifications
                    \begin{itemize}
                        \item \url{https://ieeexplore.ieee.org/document/720574/references#references}
                        \item[] (Ultimo accesso: )
                    \end{itemize}
                \end{itemize}
            }
            \subsubsection{Riferimenti informativi} {
                \begin{itemize}
                    %\item Slide del corso: T4 - Gestione di progetto
                    %\begin{itemize}
                    %    \item \url{https://www.math.unipd.it/~tullio/IS-1/2025/Dispense/T04.pdf}
                    %    \item[] (Ultimo accesso: 2025/11/07)
                    %\end{itemize}
                    \item Slide del corso: T5 - Analisi dei requisiti
                    \begin{itemize}
                        \item \url{https://www.math.unipd.it/~tullio/IS-1/2025/Dispense/T05.pdf}
                        \item[] (Ultimo accesso: 2025/11/07)
                    \end{itemize}
                    \item Slide del corso: P2 - Diagrammi dei Casi d'Uso
                    \begin{itemize}
                        \item \url{https://www.math.unipd.it/~rcardin/swea/2022/Diagrammi\%20Use\%20Case.pdf}
                        \item[] (Ultimo accesso: 2025/11/07)
                    \end{itemize}
                    \item Repository Github del Prof. Cardin
                    \begin{itemize}
                        \item \url{https://github.com/rcardin/swe-imdb}
                        \item[] (Ultimo accesso: )
                    \end{itemize}
                    \item Glossario
                \end{itemize}
            }
        }
    }

    \newpage
    \section{Descrizione del prodotto} {
        \subsection{Presentazione del prodotto} {
            DIPReader ha lo scopo di supportare la gestione e la conservazione digitale dei documenti, consentendo una transizione strutturata da archivi cartacei a soluzioni cloud avanzate. Il sistema centralizza i documenti e ne assicura autenticità, integrità, leggibilità, reperibilità e sicurezza, offrendo funzionalità di acquisizione, estrazione dei dati e consultazione. Inoltre, permette di richiedere la disponibilità offline dei documenti, garantendo sempre la conformità ai requisiti normativi. 
        }
        \subsection{Funzioni del prodotto} {
            Le funzioni del prodotto sono: 
            \begin{itemize}
                \item Rappresentazione semplificata del contenuto 
                \item Ricerca dei documenti tramite metadati
                \item Visualizzazione in anteprima di alcuni formati
                \item Salvare documenti nel computer dell'Utente
                \item Funzionamento offline
            \end{itemize}
        }
        \subsection{Caratteristiche dell'utente} {
            L’applicazione è destinata a un pubblico ampio e variegato, comprendente ricercatori, professionisti, cittadini interessati o chiunque abbia bisogno di consultare documenti archiviati. Gli utenti possono avere competenze tecniche e digitali molto diverse: dai non esperti, che desiderano trovare rapidamente informazioni, ai professionisti abituati a interpretare metadati complessi. \\
            \\
            Tutti gli utenti condividono l’esigenza di effettuare ricerche efficaci e mirate, utilizzando filtri sui metadati, consultare il contenuto dei documenti e scaricare informazioni in modo affidabile. L’interfaccia dovrà quindi garantire usabilità, chiarezza dei metadati principali e supportare la consultazione e l’estrazione dei documenti in maniera intuitiva.
        }

    }
    \newpage
    \section{Casi d'uso} {
        \subsection{Scopo} {
            Lo scopo di questa sezione è spiegare singolarmente ciascun Caso d'uso individuato.  
        }
        \subsection{Attori} {
            \begin{itemize}
                \item Utente: in questo contesto è l'Attore principale, cioè colui che intergisce attivamente con il sistema e svolge le azioni indicate dal Caso d'Uso.
                \item Motore semantico: si occupa di elaborare una ricerca semantica fatta dall'utente per produrre un risultato. In questo contesto è un Attore secondario, in quanto interagisce passivamente con il sistema.
                \item Gestore di stampa SO: si occupa di gestire la richiesta di stampa di un documento da parte dell'utente. In questo contesto è un Attore secondario, in quanto interagisce passivamente con il sistema.
            \end{itemize}
        }
        \subsection{Casi d'uso} {
            % Caricamento DIP
            \begin{usecase}
                \code{\loaddip}
                \titleUC{Caricamento DIP}
                \actor{Utente}
                \pre{
                    \item Nessun DIP è caricato nel sistema
                }
                \post{
                    \item Il DIP è caricato nel sistema
                }
                \main{
                    \item L'utente seleziona la funzione di caricamento
                    \item L'utente seleziona il pacchetto (\textbf{UC\loaddip.1})
                    \item L'utente visualizza la conferma dell'avvenuto caricamento
                }
                \alt{
                    \item[2a.] L'utente cancella l'operazione
                    \item[3a.] L'utente inserisce un pacchetto con struttura non conforme (\textbf{UC\errloaddip})
                }
                \inc{\ucref{uc:\loaddip.1}}
                \ext{\ucref{uc:\errloaddip}}
            \end{usecase}
            \begin{figure}[H]
                    \centering
                    \includegraphics[width=0.6\textwidth]{Use_Case_Diagrams/UC01.png}
                    \caption{UC01 - Caricamento DIP}
            \end{figure}
            
            \begin{usecase}
                \code{\loaddip.1}
                \titleUC{Selezione pacchetto}
                \actor{Utente}
                \pre{
                    \item L'utente ha selezionato la funzione di caricamento
                }
                \post{
                    \item L'utente ha selezionato la funzione di caricamento
                }
                \main{
                    \item Il sistema richiede il pacchetto all'utente
                    \item L'utente seleziona il pacchetto dal proprio dispositivo
                }
            \end{usecase}
            
            % Caricamento DIP con DIP già caricato
            \begin{usecase}
                \code{\remloaddip}
                \titleUC{Caricamento DIP con DIP già presente}
                \actor{Utente}
                \pre{
                    \item Un DIP è caricato nel sistema
                }
                \post{
                    \item Il DIP già caricato è rimosso dal sistema
                    \item Il nuovo DIP è caricato nel sistema
                }
                \main{
                    \item L'utente seleziona la funzione di caricamento
                    \item L'utente vede un messaggio informativo sulla necessità di rimozione del DIP presente 
                    \item L'utente rimuove il DIP presente nel sistema (\textbf{UC\removedip})
                    \item L'utente seleziona il pacchetto (\textbf{UC\loaddip.1})
                    \item L'utente visualizza la conferma dell'avvenuto caricamento
                }
                \alt{
                    \item[4a.] L'utente cancella l'operazione
                    \item[5a.] L'utente inserisce un pacchetto con struttura non conforme (\textbf{UC\errloaddip})
                }
                \inc{\ucref{uc:\loaddip.1}, \ucref{uc:\removedip}}
                \ext{\ucref{uc:\errloaddip}}
                \gen{\ucref{uc:\loaddip}}
            \end{usecase}

            % Errore struttura DIP
            \begin{usecase}
                \code{\errloaddip}
                \titleUC{Errore struttura DIP}
                \actor{Utente}
                \pre{
                    \item L'utente ha caricato un pacchetto che non rappresenta un DIP
                }
                \post{
                    \item Il DIP non viene caricato nel sistema
                    \item L'utente vede un messaggio di errore
                }
                \main{
                    \item L'utente vede un messaggio di errore relativo alla struttura non conforme del DIP
                }
            \end{usecase}
         
            % Rimozione DIP
            \begin{usecase}
                \code{\removedip}
                \titleUC{Rimozione DIP}
                \actor{Utente}
                \pre{
                    \item L'utente ha caricato un DIP nel sistema
                }
                \post{
                    \item Il DIP caricato è rimosso dal sistema
                }
                \main{
                    \item L'utente seleziona la funzionalità di chiusura
                    \item Il sistema richiede la conferma
                    \item L'utente conferma la chiusura del DIP
                }
                \alt{
                    \item[3a.] L'utente cancella l'operazione
                }
            \end{usecase}
            
            % Visualizzazione DIP
            \begin{usecase}
                \code{\viewdip}
                \titleUC{Visualizzazione contenuto DIP}
                \actor{Utente}
                \pre{
                    \item L'utente ha caricato un DIP nel sistema
                }
                \post{
                    \item L'utente vede le informazioni relative al DIP
                    \item L'utente vede i documenti contenuti nel DIP
                }
                \main{
                    \item L'utente visualizza le informazioni relative al DIP (\textbf{UC\viewdip.1})
                    \item L'utente visualizza la lista dei documenti contenuti nel DIP (\textbf{UC\viewdip.2})
                }
                \inc{\ucref{uc:\viewdip.1}, \ucref{uc:\viewdip.2}}
            \end{usecase}
            \begin{figure}[H]
                    \centering
                    \includegraphics[width=0.6\textwidth]{Use_Case_Diagrams/UC05.png}
                    \caption{UC05 - Visualizzazione contenuto DIP}
            \end{figure}
            \begin{usecase}
                \code{\viewdip.1}
                \titleUC{Visualizzazione informazioni DIP}
                \actor{Utente}
                \pre{
                    \item L'utente sta visualizzando il contenuto del DIP
                }
                \post{
                    \item L'utente vede le informazioni relative al DIP
                }
                \main{
                    \item L'utente visualizza
                    \begin{itemize}
                        \item La data di creazione del DIP
                        \item Le informazioni del programma usato per la creazione del DIP, quali
                        \begin{itemize}
                            \item Nome
                            \item Versione
                            \item Produttore
                            \item Il numero di documenti contenuti
                        \end{itemize}
                    \end{itemize}
                }
            \end{usecase}
            \begin{usecase}
                \code{\viewdip.2}
                \titleUC{Visualizzazione lista documenti contenuti}
                \actor{Utente}
                \pre{
                    \item L'utente sta visualizzando il contenuto del DIP
                }
                \post{
                    \item L'utente vede la lista di documenti contenuti nel DIP
                }
                \main{
                    \item L'utente visualizza i documenti contenuti nel DIP
                }
                \inc{\ucref{uc:\viewdip.2.1}}
            \end{usecase}
            \begin{usecase}
                \code{\viewdip.2.1}
                \titleUC{Visualizzazione singolo documento}
                \actor{Utente}
                \pre{
                    \item L'utente sta visualizzando la lista dei documenti contenuti nel DIP
                }
                \post{
                    \item L'utente vede le informazioni principali di un documento contenuto nel DIP
                }
                \main{
                    \item L'utente visualizza
                    \begin{itemize}
                        \item Il nome del file
                        \item Il suo formato
                        \item Identificativo documentale
                        \item Data in cui il documento è stato registrato
                    \end{itemize}
                }
                \inc{\ucref{uc:\viewdetdoc}}
            \end{usecase}
            \begin{usecase}
                \code{\viewdip.2.2}
                \titleUC{Visualizzazione singolo documento con allegati}
                \actor{Utente}
                \pre{
                    \item L'utente sta visualizzando la lista dei documenti contenuti nel DIP
                }
                \post{
                    \item L'utente vede le informazioni principali di un documento contenuto nel DIP
                    \item L'utente vede la lista degli allegati del documento
                }
                \main{
                    \item L'utente visualizza
                    \begin{itemize}
                        \item Il nome del file
                        \item Il suo formato
                        \item Identificativo documentale
                        \item Data in cui il documento è stato registrato
                        \item La lista dei suoi allegati (\textbf{UC\viewdip.2.2.1})
                    \end{itemize}
                }
                \inc{\ucref{uc:\viewdip.2.2.1}}
                \gen{\ucref{uc:\viewdip.2.1}}
            \end{usecase}
            \begin{usecase}
                \code{\viewdip.2.2.1}
                \titleUC{Visualizzazione lista allegati documento}
                \actor{Utente}
                \pre{
                    \item L'utente sta visualizzando un documento contenuto nel DIP
                }
                \post{
                    \item L'utente vede la lista degli allegati di un documento contenuto nel DIP
                }
                \main{
                    \item L'utente visualizza gli allegati di un documento contenuto nel DIP
                }
                \inc{\ucref{uc:\viewdip.2.2.1.1}}
            \end{usecase}
            \begin{usecase}
                \code{\viewdip.2.2.1.1}
                \titleUC{Visualizzazione singolo allegato}
                \actor{Utente}
                \pre{
                    \item L'utente sta visualizzando la lista degli allegati di un documento contenuto nel DIP
                }
                \post{
                    \item L'utente vede le informazioni principali di un allegato di un documento contenuto nel DIP
                }
                \main{
                    \item L'utente visualizza
                    \begin{itemize}
                        \item Il nome dell'allegato
                        \item Il suo formato
                        \item Identificativo allegato
                    \end{itemize}
                }
            \end{usecase}

            % Visualizzazione dettagli documento
            \begin{usecase}
                \code{\viewdetdoc}
                \titleUC{Visualizzazione dettagli documento}
                \actor{Utente}
                \pre{
                    \item L'utente sta visualizzando un singolo documento
                }
                \post{
                    \item L'utente vede i dettagli di un documento
                }
                \main{
                    \item L'utente seleziona la funzionalità di visualizzazione dei dettagli di un documento
                    \item L'utente visualizza i dettagli di un documento
                }
            \end{usecase}
            
            % Visualizzazione delle informazioni di conservazione
            \begin{usecase}
                \code{\viewconsinfo}
                \titleUC{Visualizzazione informazioni di conservazione}
                \actor{Utente}
                \pre{
                    \item L'utente ha caricato un DIP
                    \item L'utente sta visualizzando una lista di documenti contenuti nel DIP
                }
                \post{
                    \item L’utente visualizza i dettagli relativi alla validità legale e allo stato del processo di conservazione dei documenti
                }
                \main{
                    \item L'utente visualizza degli elementi grafici che gli indicano la validità e lo stato del processo di conservazione dei documenti
                }
            \end{usecase}
            \begin{figure}[H]
                    \centering
                    \includegraphics[width=0.5\textwidth]{Use_Case_Diagrams/UC07.png}
                    \caption{UC07 - Visualizzazione informazioni di conservazione}
            \end{figure}

            % Criteri di ricerca (metadati)
            \begin{usecase}
                \code{\mdparamImpCrit}
                \titleUC{Inserimento criterio di ricerca "Impronta crittografica"}
                \actor{Utente}
                \pre{
                    \item L'utente ha caricato un DIP
                    \item L'utente sta visualizzando una lista di documenti contenuti nel DIP
                }
                \post{
                    \item L'utente ha configurato un criterio di ricerca
                }
                \main{
                    \item L'utente seleziona la funzionalità di aggiunta criterio di ricerca
                    \item L'utente seleziona il criterio "Impronta crittografica"
                    \item L'utente inserisce l'hash del documento (\textbf{UC\mdparamImpCrit.1})
                    \item L'utente seleziona l'algoritmo applicato (\textbf{UC\mdparamImpCrit.2})
                }
                \inc{\ucref{uc:\mdparamImpCrit.1}, \ucref{uc:\mdparamImpCrit.2}}
            \end{usecase}
            \begin{figure}[H]
                    \centering
                    \includegraphics[width=0.6\textwidth]{Use_Case_Diagrams/UC08.png}
                    \caption{UC08 - Inserimento criterio di ricerca ”Impronta crittografica”}
            \end{figure}
            \begin{usecase}
                \code{\mdparamImpCrit.1}
                \titleUC{Inserimento hash del documento}
                \actor{Utente}
                \pre{
                    \item L'utente sta configurando il criterio di ricerca "Impronta crittografica"
                }
                \post{
                    \item L'utente ha inserito l'hash del documento
                }
                \main{
                    \item L'utente inserisce l'hash del documento
                }
            \end{usecase}
            \begin{usecase}
                \code{\mdparamImpCrit.2}
                \titleUC{Selezione algoritmo di hashing}
                \actor{Utente}
                \pre{
                    \item L'utente sta configurando il criterio di ricerca "Impronta crittografica"
                }
                \post{
                    \item L'utente ha selezionato l'algoritmo di hashing
                }
                \main{
                    \item L'utente seleziona l'algoritmo di hashing applicato tra quelli disponibili
                }
            \end{usecase}

            \begin{usecase}
                \code{\mdparamIdDoc}
                \titleUC{Inserimento criterio di ricerca "Identificativo documento"}
                \actor{Utente}
                \pre{
                    \item L'utente ha caricato un DIP
                    \item L'utente sta visualizzando una lista di documenti contenuti nel DIP
                }
                \post{
                    \item L'utente ha configurato un criterio di ricerca
                }
                \main{
                    \item L'utente seleziona la funzionalità di aggiunta criterio di ricerca
                    \item L'utente seleziona il criterio "Identificativo documento"
                    \item L'utente inserisce il valore
                }
            \end{usecase}
            \begin{figure}[H]
                    \centering
                    \includegraphics[width=0.5\textwidth]{Use_Case_Diagrams/UC09.png}
                    \caption{UC09 - Inserimento criterio di ricerca ”Identificativo documento”}
            \end{figure}

            \begin{usecase}
                \code{\mdparamModFor}
                \titleUC{Inserimento criterio di ricerca "Modalità di formazione"}
                \actor{Utente}
                \pre{
                    \item L'utente ha caricato un DIP
                    \item L'utente sta visualizzando una lista di documenti contenuti nel DIP
                }
                \post{
                    \item L'utente ha configurato un criterio di ricerca
                }
                \main{
                    \item L'utente seleziona la funzionalità di aggiunta criterio di ricerca
                    \item L'utente seleziona il criterio "Modalità di formazione"
                    \item L'utente seleziona i valori tra le opzioni
                    \begin{itemize}
                        \item "Creazione tramite l'utilizzo di strumenti software che assicurino la produzione di documenti nei formati previsti nell'Allegato 2 delle Linee Guida"
                        \item "Acquisizione di un documento informatico per via telematica o su supporto informatico, acquisizione della copia per immagine su supporto informatico di un documento analogico, acquisizione della copia informatica di un documento analogico"
                        \item "Memorizzazione su supporto informatico in formato digitale delle informazioni risultanti da transazioni o processi informatici o dalla presentazione telematica di dati attraverso moduli o formulari resi disponibili all'utente"
                        \item "Generazione o raggruppamento anche in via automatica di un insieme di dati o registrazioni, provenienti da una o più banche dati, anche appartenenti a più soggetti interoperanti, secondo una struttura logica predeterminata e memorizzata in forma statica"
                    \end{itemize}
                }
            \end{usecase}
            \begin{figure}[H]
                    \centering
                    \includegraphics[width=0.5\textwidth]{Use_Case_Diagrams/UC10.png}
                    \caption{UC10 - Inserimento criterio di ricerca ”Modalità di formazione”}
            \end{figure}

            \begin{usecase}
                \code{\mdparamTipDoc}
                \titleUC{Inserimento criterio di ricerca "Tipologia documentale"}
                \actor{Utente}
                \pre{
                    \item L'utente ha caricato un DIP
                    \item L'utente sta visualizzando una lista di documenti contenuti nel DIP
                }
                \post{
                    \item L'utente ha configurato un criterio di ricerca
                }
                \main{
                    \item L'utente seleziona la funzionalità di aggiunta criterio di ricerca
                    \item L'utente seleziona il criterio "Tipologia documentale"
                    \item L'utente inserisce il valore
                }
            \end{usecase}
            \begin{figure}[H]
                    \centering
                    \includegraphics[width=0.5\textwidth]{Use_Case_Diagrams/UC11.png}
                    \caption{UC11 - Inserimento criterio di ricerca ”Tipologia documentale”}
            \end{figure}

            \begin{usecase}
                \code{\mdparamTipFlu}
                \titleUC{Inserimento criterio di ricerca "Tipologia di flusso"}
                \actor{Utente}
                \pre{
                    \item L'utente ha caricato un DIP
                    \item L'utente sta visualizzando una lista di documenti contenuti nel DIP
                }
                \post{
                    \item L'utente ha configurato un criterio di ricerca
                }
                \main{
                    \item L'utente seleziona la funzionalità di aggiunta criterio di ricerca
                    \item L'utente seleziona il criterio "Tipologia di flusso"
                    \item L'utente seleziona i valori tra le opzioni
                    \begin{itemize}
                        \item "U (In uscita)"
                        \item "E (In entrata)"
                        \item "I (Interno)"
                    \end{itemize}
                }
            \end{usecase}
            \begin{figure}[H]
                    \centering
                    \includegraphics[width=0.5\textwidth]{Use_Case_Diagrams/UC12.png}
                    \caption{UC12 - Inserimento criterio di ricerca ”Tipologia di flusso”}
            \end{figure}

            \begin{usecase}
                \code{\mdparamTipReg}
                \titleUC{Inserimento criterio di ricerca "Tipo registro"}
                \actor{Utente}
                \pre{
                    \item L'utente ha caricato un DIP
                    \item L'utente sta visualizzando una lista di documenti contenuti nel DIP
                }
                \post{
                    \item L'utente ha configurato un criterio di ricerca
                }
                \main{
                    \item L'utente seleziona la funzionalità di aggiunta criterio di ricerca
                    \item L'utente seleziona il criterio "Tipo registro"
                    \item L'utente seleziona il valore tra le opzioni
                    \begin{itemize}
                        \item "Nessuno"
                        \item "Protocollo ordinario/Protocollo emergenza"
                        \item "Repertorio/Registro"
                    \end{itemize}
                }
            \end{usecase}
            \begin{figure}[H]
                    \centering
                    \includegraphics[width=0.5\textwidth]{Use_Case_Diagrams/UC13.png}
                    \caption{UC13 - Inserimento criterio di ricerca ”Tipo registro”}
            \end{figure}

            \begin{usecase}
                \code{\mdparamDatReg}
                \titleUC{Inserimento criterio di ricerca "Data registrazione"}
                \actor{Utente}
                \pre{
                    \item L'utente ha caricato un DIP
                    \item L'utente sta visualizzando una lista di documenti contenuti nel DIP
                }
                \post{
                    \item L'utente ha configurato un criterio di ricerca
                }
                \main{
                    \item L'utente seleziona la funzionalità di aggiunta criterio di ricerca
                    \item L'utente seleziona il criterio "Data registrazione"
                    \item L'utente inserisce il valore
                }
            \end{usecase}
            \begin{figure}[H]
                    \centering
                    \includegraphics[width=0.5\textwidth]{Use_Case_Diagrams/UC14.png}
                    \caption{UC14 - Inserimento criterio di ricerca ”Data registrazione”}
            \end{figure}

            \begin{usecase}
                \code{\mdparamOraReg}
                \titleUC{Inserimento criterio di ricerca "Ora registrazione"}
                \actor{Utente}
                \pre{
                    \item L'utente ha caricato un DIP
                    \item L'utente sta visualizzando una lista di documenti contenuti nel DIP
                }
                \post{
                    \item L'utente ha configurato un criterio di ricerca
                }
                \main{
                    \item L'utente seleziona la funzionalità di aggiunta criterio di ricerca
                    \item L'utente seleziona il criterio "Ora registrazione"
                    \item L'utente inserisce il valore
                }
            \end{usecase}
            \begin{figure}[H]
                    \centering
                    \includegraphics[width=0.5\textwidth]{Use_Case_Diagrams/UC15.png}
                    \caption{UC15 - Inserimento criterio di ricerca ”Ora registrazione”}
            \end{figure}

            \begin{usecase}
                \code{\mdparamNumDoc}
                \titleUC{Inserimento criterio di ricerca "Numero documento"}
                \actor{Utente}
                \pre{
                    \item L'utente ha caricato un DIP
                    \item L'utente sta visualizzando una lista di documenti contenuti nel DIP
                }
                \post{
                    \item L'utente ha configurato un criterio di ricerca
                }
                \main{
                    \item L'utente seleziona la funzionalità di aggiunta criterio di ricerca
                    \item L'utente seleziona il criterio "Numero documento"
                    \item L'utente inserisce il valore
                }
            \end{usecase}
            \begin{figure}[H]
                    \centering
                    \includegraphics[width=0.5\textwidth]{Use_Case_Diagrams/UC16.png}
                    \caption{UC16 - Inserimento criterio di ricerca ”Numero documento”}
            \end{figure}

            \begin{usecase}
                \code{\mdparamCodReg}
                \titleUC{Inserimento criterio di ricerca "Codice registro"}
                \actor{Utente}
                \pre{
                    \item L'utente ha caricato un DIP
                    \item L'utente sta visualizzando una lista di documenti contenuti nel DIP
                }
                \post{
                    \item L'utente ha configurato un criterio di ricerca
                }
                \main{
                    \item L'utente seleziona la funzionalità di aggiunta criterio di ricerca
                    \item L'utente seleziona il criterio "Codice registro"
                    \item L'utente inserisce il valore
                }
            \end{usecase}
            \begin{figure}[H]
                    \centering
                    \includegraphics[width=0.5\textwidth]{Use_Case_Diagrams/UC17.png}
                    \caption{UC17 - Inserimento criterio di ricerca ”Codice registro”}
            \end{figure}
        }
    }
    \newpage
    \section{Requisiti} {
       
        \subsection{Requisiti Funzionali} {
            
        }
        \subsection{Requisiti di Qualità} {
           
        }
        \subsection{Requisiti di Vincolo} {
           
        }
    }

\end{document}